% OpenStax Calculus Volume 2 -- Section 5.5 Exercises: Alternating Series
% Exercises reproduced from OpenStax, licensed under CC BY 4.0.
% Source: https://openstax.org/books/calculus-volume-2/pages/5-5-alternating-series
\documentclass{exercises}
\title{OpenStax Calculus Volume 2}
\subtitle{Section 5.5 Exercises: Alternating Series}
\sourceurl{https://openstax.org/books/calculus-volume-2/pages/5-5-alternating-series}
\author{}
\date{}

\begin{document}
\maketitle


State whether each of the following series converges absolutely, conditionally, or not at all.

\begin{enumerate}[start=1]
\item $\sum_{n=1}^{\infty}{(-1)}^{n+1}\frac{n}{n+3}$
\item $\sum_{n=1}^{\infty}{(-1)}^{n+1}\frac{\sqrt{n}+1}{\sqrt{n}+3}$
\item $\sum_{n=1}^{\infty}{(-1)}^{n+1}\frac{1}{\sqrt{n+3}}$
\item $\sum_{n=1}^{\infty}{(-1)}^{n+1}\frac{\sqrt{n+3}}{n}$
\item $\sum_{n=1}^{\infty}{(-1)}^{n+1}\frac{1}{n!}$
\item $\sum_{n=1}^{\infty}{(-1)}^{n+1}\frac{3^{n}}{n!}$
\item $\sum_{n=1}^{\infty}{(-1)}^{n+1}{(\frac{n-1}{n})}^{n}$
\item $\sum_{n=1}^{\infty}{(-1)}^{n+1}{(\frac{n+1}{n})}^{n}$
\item $\sum_{n=1}^{\infty}{(-1)}^{n+1}\sin^{2}n$
\item $\sum_{n=1}^{\infty}{(-1)}^{n+1}\cos^{2}n$
\item $\sum_{n=1}^{\infty}{(-1)}^{n+1}\sin^{2}(1/n)$
\item $\sum_{n=1}^{\infty}{(-1)}^{n+1}\cos^{2}(1/n)$
\item $\sum_{n=1}^{\infty}{(-1)}^{n+1}\ln(1/n)$
\item $\sum_{n=1}^{\infty}{(-1)}^{n+1}\ln(1+\frac{1}{n})$
\item $\sum_{n=1}^{\infty}{(-1)}^{n+1}\frac{n^{2}}{1+n^{4}}$
\item $\sum_{n=1}^{\infty}{(-1)}^{n+1}\frac{n^{e}}{1+n^{\pi}}$
\item $\sum_{n=1}^{\infty}{(-1)}^{n+1}2^{1/n}$
\item $\sum_{n=1}^{\infty}{(-1)}^{n+1}n^{1/n}$
\item $\sum_{n=1}^{\infty}{(-1)}^{n}(1-n^{1/n})$ (Hint: $n^{1/n}\approx1+\ln(n)/n$ for large $n$.)
\item $\sum_{n=1}^{\infty}{(-1)}^{n+1}n(1-\cos(\frac{1}{n}))$ (Hint: $\cos(1/n)\approx1-1/n^{2}$ for large $n$.)
\item $\sum_{n=1}^{\infty}{(-1)}^{n+1}(\sqrt{n+1}-\sqrt{n})$ (Hint: Rationalize the numerator.)
\item $\sum_{n=1}^{\infty}{(-1)}^{n+1}(\frac{1}{\sqrt{n}}-\frac{1}{\sqrt{n+1}})$ (Hint: Find common denominator then rationalize numerator.)
\item $\sum_{n=1}^{\infty}{(-1)}^{n+1}(\ln(n+1)-\ln n)$
\item $\sum_{n=1}^{\infty}{(-1)}^{n+1}n(\tan^{-1}(n+1)-\tan^{-1}n)$ (Hint: Use the Mean Value Theorem.)
\item $\sum_{n=1}^{\infty}{(-1)}^{n+1}({(n+1)}^{2}-n^{2})$
\item $\sum_{n=1}^{\infty}{(-1)}^{n+1}(\frac{1}{n}-\frac{1}{n+1})$
\item $\sum_{n=1}^{\infty}\frac{\cos(n\pi)}{n}$
\item $\sum_{n=1}^{\infty}\frac{\cos(n\pi)}{n^{1/n}}$
\item $\sum_{n=1}^{\infty}\frac{1}{n}\sin(\frac{n\pi}{2})$
\item $\sum_{n=1}^{\infty}\sin(n\pi/2)\sin(1/n)$
\end{enumerate}

In each of the following problems, use the estimate $|R_{N}|\le b_{N+1}$ to find a value of $N$ that guarantees that the sum of the first $N$ terms of the alternating series $\sum_{n=1}^{\infty}{(-1)}^{n+1}b_{n}$ differs from the infinite sum by at most the given error. Calculate the partial sum $S_{N}$ for this $N$.

\begin{enumerate}[resume]
\item \techrequired $b_{n}=1/n$, error $<10^{-5}$
\item \techrequired $b_{n}=1/\ln(n)$, $n\ge2$, error $<10^{-1}$
\item \techrequired $b_{n}=1/\sqrt{n}$, error $<10^{-3}$
\item \techrequired $b_{n}=1/2^{n}$, error $<10^{-6}$
\item \techrequired $b_{n}=\ln(1+\frac{1}{n})$, error $<10^{-3}$
\item \techrequired $b_{n}=1/n^{2}$, error $<10^{-6}$
\end{enumerate}

For the following exercises, indicate whether each of the following statements is true or false. If the statement is false, provide an example in which it is false.

\begin{enumerate}[resume]
\item If $b_{n}\ge0$ is decreasing and $\lim_{n\to \infty}b_{n}=0$, then $\sum_{n=1}^{\infty}(b_{2n-1}-b_{2n})$ converges absolutely.
\item If $b_{n}\ge0$ is decreasing, then $\sum_{n=1}^{\infty}(b_{2n-1}-b_{2n})$ converges absolutely.
\item If $b_{n}\ge0$ and $\lim_{n\to \infty}b_{n}=0$, then $\sum_{n=1}^{\infty}(\frac{1}{2}(b_{3n-2}+b_{3n-1})-b_{3n})$ converges.
\item If $b_{n}\ge0$ is decreasing and $\sum_{n=1}^{\infty}(b_{3n-2}+b_{3n-1}-b_{3n})$ converges, then $\sum_{n=1}^{\infty}b_{3n-2}$ converges.
\item If $b_{n}\ge0$ is decreasing and $\sum_{n=1}^{\infty}{(-1)}^{n-1}b_{n}$ converges conditionally but not absolutely, then $b_{n}$ does not tend to zero.
\item Let $a_{n}^{+}=a_{n}$ if $a_{n}\ge0$ and $a_{n}^{-}=-a_{n}$ if $a_{n}<0$. (Also, $a_{n}^{+}=0$ if $a_{n}<0$ and $a_{n}^{-}=0$ if $a_{n}\ge0$.) If $\sum_{n=1}^{\infty}a_{n}$ converges conditionally but not absolutely, then neither $\sum_{n=1}^{\infty}a_{n}^{+}$ nor $\sum_{n=1}^{\infty}a_{n}^{-}$ converges.
\item Suppose that $a_{n}$ is a sequence of positive real numbers and that $\sum_{n=1}^{\infty}a_{n}$ converges. Suppose that $b_{n}$ is an arbitrary sequence of ones and minus ones. Does $\sum_{n=1}^{\infty}a_{n}b_{n}$ necessarily converge?
\item Suppose that $a_{n}$ is a sequence such that $\sum_{n=1}^{\infty}a_{n}b_{n}$ converges for every possible sequence $b_{n}$ of zeros and ones. Does $\sum_{n=1}^{\infty}a_{n}$ converge absolutely?
\end{enumerate}

The following series do not satisfy the hypotheses of the alternating series test as stated.


In each case, state which hypothesis is not satisfied. State whether the series converges absolutely.

\begin{enumerate}[resume]
\item $\sum_{n=1}^{\infty}{(-1)}^{n+1}\frac{\sin^{2}n}{n}$
\item $\sum_{n=1}^{\infty}{(-1)}^{n+1}\frac{\cos^{2}n}{n}$
\item $1+\frac{1}{2}-\frac{1}{3}-\frac{1}{4}+\frac{1}{5}+\frac{1}{6}-\frac{1}{7}-\frac{1}{8}+\cdots$
\item $1+\frac{1}{2}-\frac{1}{3}+\frac{1}{4}+\frac{1}{5}-\frac{1}{6}+\frac{1}{7}+\frac{1}{8}-\frac{1}{9}+\cdots$
\item Show that the alternating series $1-\frac{1}{2}+\frac{1}{2}-\frac{1}{4}+\frac{1}{3}-\frac{1}{6}+\frac{1}{4}-\frac{1}{8}+\cdots$ does not converge. What hypothesis of the alternating series test is not met?
\item Suppose that $\sum a_{n}$ converges absolutely. Show that the series consisting of the positive terms $a_{n}$ also converges.
\item Show that the alternating series $\frac{2}{3}-\frac{3}{5}+\frac{4}{7}-\frac{5}{9}+\cdots$ does not converge. What hypothesis of the alternating series test is not met?
\item The formula $\cos \theta=1-\frac{\theta^{2}}{2!}+\frac{\theta^{4}}{4!}-\frac{\theta^{6}}{6!}+\cdots$ will be derived in the next chapter. Use the remainder $|R_{N}|\le b_{N+1}$ to find a bound for the error in estimating $\cos \theta$ by the fifth partial sum $1-\theta^{2}/2!+\theta^{4}/4!-\theta^{6}/6!+\theta^{8}/8!$ for $\theta=1$, $\theta=\pi/6$, and $\theta=\pi$.
\item The formula $\sin \theta=\theta-\frac{\theta^{3}}{3!}+\frac{\theta^{5}}{5!}-\frac{\theta^{7}}{7!}+\cdots$ will be derived in the next chapter. Use the remainder $|R_{N}|\le b_{N+1}$ to find a bound for the error in estimating $\sin \theta$ by the fifth partial sum $\theta-\theta^{3}/3!+\theta^{5}/5!-\theta^{7}/7!+\theta^{9}/9!$ for $\theta=1$, $\theta=\pi/6$, and $\theta=\pi$.
\item How many terms in $\cos \theta=1-\frac{\theta^{2}}{2!}+\frac{\theta^{4}}{4!}-\frac{\theta^{6}}{6!}+\cdots$ are needed to approximate $\cos 1$ accurate to an error of at most $0.00001$?
\item How many terms in $\sin \theta=\theta-\frac{\theta^{3}}{3!}+\frac{\theta^{5}}{5!}-\frac{\theta^{7}}{7!}+\cdots$ are needed to approximate $\sin 1$ accurate to an error of at most $0.00001$?
\item Sometimes the alternating series $\sum_{n=1}^{\infty}{(-1)}^{n-1}b_{n}$ converges to a certain fraction of an absolutely convergent series $\sum_{n=1}^{\infty}b_{n}$ at a faster rate. Given that $\sum_{n=1}^{\infty}\frac{1}{n^{2}}=\frac{\pi^{2}}{6}$, find $\frac{\pi^{2}}{12}=1-\frac{1}{2^{2}}+\frac{1}{3^{2}}-\frac{1}{4^{2}}+\cdots$. Which of the series $6\sum_{n=1}^{\infty}\frac{1}{n^{2}}$ and $12\sum_{n=1}^{\infty}\frac{{(-1)}^{n-1}}{n^{2}}$ gives a better estimation of $\pi^{2}$ using $1000$ terms?
\end{enumerate}

The following alternating series converge to given multiples of $\pi$. Find the value of $N$ predicted by the remainder estimate such that the $N\text{th}$ partial sum of the series accurately approximates the left-hand side to within the given error. Find the minimum $N$ for which the error bound holds, and give the desired approximate value in each case. Up to $15$ decimal places, $\pi=3.141592653589793\ldots$.

\begin{enumerate}[resume]
\item \techrequired $\frac{\pi}{4}=\sum_{n=0}^{\infty}\frac{{(-1)}^{n}}{2n+1}$, error $<0.0001$
\item \techrequired $\frac{\pi}{\sqrt{12}}=\sum_{k=0}^{\infty}\frac{{(-3)}^{-k}}{2k+1}$, error $<0.0001$
\item \techrequired The series $\sum_{n=0}^{\infty}\frac{\sin(x+\pi n)}{x+\pi n}$ plays an important role in signal processing. Show that $\sum_{n=0}^{\infty}\frac{\sin(x+\pi n)}{x+\pi n}$ converges whenever $0<x<\pi$. (Hint: Use the formula for the sine of a sum of angles.)
\item \techrequired If $\sum_{n=1}^{N}{(-1)}^{n-1}\frac{1}{n}\to \ln 2$, what is $1+\frac{1}{3}+\frac{1}{5}-\frac{1}{2}-\frac{1}{4}-\frac{1}{6}+\frac{1}{7}+\frac{1}{9}+\frac{1}{11}-\frac{1}{8}-\frac{1}{10}-\frac{1}{12}+\cdots$?
\item \techrequired Plot the partial sum $\sum_{n=1}^{100}\frac{\cos(2\pi nx)}{n}$ for $0\le x<1$. Explain why $\sum_{n=1}^{\infty}\frac{\cos(2\pi nx)}{n}$ diverges when $x=0,1$. How does the series behave for other $x$?
\item \techrequired Plot the series $\sum_{n=1}^{100}\frac{\sin(2\pi nx)}{n}$ for $0\le x<1$ and comment on its behavior
\item \techrequired Plot the series $\sum_{n=1}^{100}\frac{\cos(2\pi nx)}{n^{2}}$ for $0\le x<1$ and describe its graph.
\item \techrequired The alternating harmonic series converges because of cancellation among its terms. Its sum is known because the cancellation can be described explicitly. A random harmonic series is one of the form $\sum_{n=1}^{\infty}\frac{s_{n}}{n}$, where $s_{n}$ is a randomly generated sequence of $\pm1$'s in which the values $\pm1$ are equally likely to occur. Use a random number generator to produce $1000$ random $\pm1$'s and plot the partial sums $S_{N}=\sum_{n=1}^{N}\frac{s_{n}}{n}$ of your random harmonic sequence for $N=1$ to $1000$. Compare to a plot of the first $1000$ partial sums of the harmonic series.
\item \techrequired Estimates of $\sum_{n=1}^{\infty}\frac{1}{n^{2}}$ can be accelerated by writing its partial sums as $\sum_{n=1}^{N}\frac{1}{n^{2}}=\sum_{n=1}^{N}\frac{1}{n(n+1)}+\sum_{n=1}^{N}\frac{1}{n^{2}(n+1)}$ and recalling that $\sum_{n=1}^{N}\frac{1}{n(n+1)}=1-\frac{1}{N+1}$ converges to one as $N\to \infty$. Compare the estimate of $\pi^{2}/6$ using the sums $\sum_{n=1}^{1000}\frac{1}{n^{2}}$ with the estimate using $1+\sum_{n=1}^{1000}\frac{1}{n^{2}(n+1)}$.
\item \techrequired The Euler transform rewrites $S=\sum_{n=0}^{\infty}{(-1)}^{n}b_{n}$ as $S=\sum_{n=0}^{\infty}2^{-n-1}\sum_{m=0}^{n}{(-1)}^{m}\binom{n}{m}b_{m}$. For the alternating harmonic series, it takes the form $\ln(2)=\sum_{n=1}^{\infty}\frac{{(-1)}^{n-1}}{n}=\sum_{n=1}^{\infty}\frac{1}{n2^{n}}$. Compute partial sums of $\sum_{n=1}^{\infty}\frac{1}{n2^{n}}$ until they approximate $\ln(2)$ accurate to within $0.0001$. How many terms are needed? Compare this answer to the number of terms of the alternating harmonic series needed to estimate $\ln(2)$.
\item \techrequired In the text it was stated that a conditionally convergent series can be rearranged to converge to any number. Here is a slightly simpler, but similar, fact. If $a_{n}\ge0$ is such that $a_{n}\to0$ as $n\to \infty$ but $\sum_{n=1}^{\infty}a_{n}$ diverges, then, given any number $A$ there is a sequence $s_{n}$ of $\pm1\text{'s}$ such that $\sum_{n=1}^{\infty}a_{n}s_{n}\to A$. Show this for $A>0$ as follows.
\begin{enumerate}[label=(\alph*)]
  \item Recursively define $s_{n}$ by $s_{n}=1$ if $S_{n-1}=\sum_{k=1}^{n-1}a_{k}s_{k}<A$ and $s_{n}=-1$ otherwise.
  \item Explain why eventually $S_{n}\ge A$, and for any $m$ larger than this $n$, $A-a_{m}\le S_{m}\le A+a_{m}$.
  \item Explain why this implies that $S_{n}\to A$ as $n\to \infty$.
\end{enumerate}
\end{enumerate}

\end{document}
